% Главные размерения и коэффициенты

\nomenclature[D]{$L$}{Длина судна}

\nomenclature[D]{$L_{\text{НБ}}$}{Длина судна наибольшая}

\nomenclature[D]{$L_{\text{КВЛ}}$}{Длина судна по КВЛ}

\nomenclature[D]{$L_{\text{ПП}}$}{Длина судна между перпендикулярами}

\nomenclature[D]{$B$}{Ширина судна}

\nomenclature[D]{$B_{\text{НБ}}$}{Ширина судна наибольшая}

\nomenclature[D]{$B_{\text{КВЛ}}$}{Ширина судна по КВЛ}

\nomenclature[D]{$H$}{Высота борта}

\nomenclature[D]{$T$}{Осадка}

\nomenclature[D]{$T_{\text{КВЛ}}$}{Осадка по КВЛ}

\nomenclature[D]{$T_{\text{Н}}$}{Осадка носом (на носовом перпендикуляре)}

\nomenclature[D]{$T_{\text{К}}$}{Осадка кормой (на кормовом перпендикуляре)}

\nomenclature[D]{$A_L$}{Площадь погруженной части диаметральной плоскости}

\nomenclature[D]{$\Omega$}{Площадь шпангоута}

\nomenclature[D]{$x_g$}{Абсцисса центра тяжести судна}

\nomenclature[D]{$V_{\text{П}}$}{Водоизмещение по верхнюю палубу}

\nomenclature[D]{$v$}{Объем отсека}

\nomenclature[D]{$v_{\text{в}}$}{Объем отсека с учетом коэффициента проницаемости}

\nomenclature[D]{$V_X$}{Интегральная кривая  строевой по шпангоутам}

\nomenclature[D]{$\mu$}{Коэффициент проницаемости}

\nomenclature[D]{$L/B$}{Отношение длины к ширине}

\nomenclature[D]{$B/T$}{Отношение ширины к осадке}

\nomenclature[D]{$B/H$}{Отношение ширины к высоте борта}

\nomenclature[D]{$\delta$}{Коэффициент общей полноты}

\nomenclature[D]{$\beta$}{Коэффициент полноты площади шпангоута}

\nomenclature[D]{$\alpha$}{Коэффициент полноты ватерлинии}

\nomenclature[D]{$\phi$}{Коэффициент продольной полноты}

\nomenclature[D]{$\chi$}{Коэффициент вертикальной полноты}

\nomenclature[D]{$\sigma$}{Коэффициент полноты диаметральной плоскости}

% Элементы ТЧ и остойчивость

\nomenclature[T]{$V$}{Объемное водоизмещение}

\nomenclature[T]{$x_c$}{Абсцисса центра величины}

\nomenclature[T]{$z_c$}{Аппликата центра величины}

\nomenclature[T]{$S$}{Площадь ватерлинии}

\nomenclature[T]{$x_f$}{Абсцисса центра тяжести площади ватерлинии}

\nomenclature[T]{$I_x$}{Момент инерции площади ватерлинии относительно оси $X$}

\nomenclature[T]{$I_y$}{Момент инерции площади ватерлинии относительно оси $Y$}

\nomenclature[T]{$I_{yf}$}{Момент инерции площади ватерлинии относительно поперечной центральной оси}

\nomenclature[T]{$r$}{Поперечный метацентрический радиус}

\nomenclature[T]{$R$}{Продольный метацентрический радиус}

\nomenclature[T]{$z_m$}{Возвышение поперечного метацентра}

\nomenclature[T]{$z_M$}{Возвышение продольного метацентра}

\nomenclature[T]{$M_y$}{Статический момент площади ватерлинии относительно оси $Y$}

\nomenclature[T]{$M_{yz}$}{Статический момент объема относительно плоскости мидель-шпангоута (плоскости $YZ$)}

\nomenclature[T]{$M_{xy}$}{Статический момент объема относительно основной плоскости (плоскости $XY$)}

\nomenclature[T]{$l_{\Theta}$}{Плечо статической остойчивости}

\nomenclature[T]{$l_{\text{ф}}$}{Плечо остойчивости формы}

\nomenclature[T]{$a$}{Расстояние от центра тяжести до центра величины}

\nomenclature[T]{$d$}{Плечо динамической остойчивости}

\nomenclature[T]{$_{\Theta}$}{Индекс относится к величинам, определенным при наклонении на угол $\Theta$}

\nomenclature[T]{$y_{\Theta}$}{Ордината центра величины при наклонении на угол $\Theta$}

\nomenclature[T]{$z_{\Theta}$}{Аппликата центра величины при наклонении на угол $\Theta$}

% Ординаты

\nomenclature[O]{$\bar{y}$}{Безрамерная ордината точки пересечения шпангоута и ватерлинии}

\nomenclature[O]{$\text{П}$}{Безрамерная ордината точки пересечения шпангоута и линии борта}

\nomenclature[O]{$\bar{z}$}{Безрамерная аппликата точки пересечения шпангоута и линии борта}

\nomenclature[O]{$\bar{z}_1$}{Безрамерная аппликата точки пересечения шпангоута и первого батокса}

\nomenclature[O]{$\bar{z}_2$}{Безрамерная аппликата точки пересечения шпангоута и второго батокса}

\nomenclature[O]{$\bar{z}_{\text{Ф}}$}{Безрамерная аппликата точки пересечения форштевня с линией борта}

\nomenclature[O]{$\bar{z}_{\text{А}}$}{Безрамерная аппликата точки пересечения ахтерштевня с линией борта}

\nomenclature[O]{$\bar{x}_{\text{Ф}}$}{Безрамерная абсцисса точки пересечения форштевня и ватерлинии}

\nomenclature[O]{$\bar{x}_{\text{А}}$}{Безрамерная абсцисса точки пересечения ахтерштевня и ватерлинии}

\nomenclature[O]{$y_{\text{П}}$}{Размерная ордината точки пересечения шпангоута и линии борта}

\nomenclature[O]{$y_{\text{пр}}$}{Размерная ордината правого борта}

\nomenclature[O]{$y_{\text{лб}}$}{Размерная ордината левого борта}

\nomenclature[O]{$x_{\text{Н}}$}{Абсцисса носовой точки ватерлинии}

\nomenclature[O]{$x_{\text{К}}$}{Абсцисса кормовой точки ватерлинии}

\nomenclature[O]{$f$}{Погибь бимсов}

\nomenclature[O]{$\Delta L_{\text{Ф}}$}{Длина от точки притыкания ватерлинии к ДП до ближайшего шпангоута в носу}

\nomenclature[O]{$\Delta L_{\text{А}}$}{Длина от точки притыкания ватерлинии к ДП до ближайшего шпангоута в корме}

\nomenclature[O]{$A_{\text{Ф}}$}{Площадь треугольника ватерлинии в носу}

\nomenclature[O]{$A_{\text{А}}$}{Площадь треугольника ватерлинии в корме}

\nomenclature[O]{$K_{\text{Ф}}$}{Коэффициент учета притыкания ватерлинии к ДП в носу}

\nomenclature[O]{$K_{\text{А}}$}{Коэффициент учета притыкания ватерлинии к ДП в корме}

% Интегрирования

\nomenclature[I]{$f(x)$}{Функция, зависящая от переменной $x$}

\nomenclature[I]{$f_{\text{н}}$}{Значение функции на крайнем носовом шпангоуте}

\nomenclature[I]{$f_{\text{к}}$}{Значение функции на крайнем кормовом шпангоуте}

\nomenclature[I]{$f_{\text{т}}$}{Значение функции на транце}

\nomenclature[I]{$i$}{Номер шпангоута}

% Спуск

\nomenclature[S]{$D_{\text{С}}$}{Спусковой вес}

\nomenclature[S]{$L$}{Длина полоза}

\nomenclature[S]{$L_1$}{Длина переднего конца полоза}

\nomenclature[S]{$L_2$}{Длина заднего конца полоза}

\nomenclature[S]{$T_0$}{Расчетная глубина воды на пороге}

\nomenclature[S]{$\beta$}{Уклон спусковых дорожек}

\nomenclature[S]{$\rho$}{Плотность воды}

\nomenclature[S]{$\alpha$}{Уклон линии киля}

\nomenclature[S]{$\lambda$}{Длина подводной части спусковых дорожек}

\nomenclature[S]{$T_1$}{Осадка передним концом полоза}

\nomenclature[S]{$T_2$}{Осадка задним концом полоза}

\nomenclature[S]{$N_{\text{Б}}$}{Баксово давление}

\nomenclature[S]{$W$}{Объемное водоизмещение погруженной части судна}

\nomenclature[S]{$M_D$}{Момент силы веса}

\nomenclature[S]{$M_W$}{Момент силы плавучести}

% Сокращения

\nomenclature[A]{CAD}{Автоматизированное проектирование (Computer-aided design)}

\nomenclature[A]{ВЛ}{Ватерлиния}

\nomenclature[A]{ДП}{Диаметральная плоскость}

\nomenclature[A]{ЕСКД}{Единая система конструкторской документации}

\nomenclature[A]{ОП}{Основная плоскость}

\nomenclature[A]{САПР}{Система автоматизированного проектирования}


\printnomenclature
